% Versión en español del perfil de AI Engineer.
%
% ALCANCE (ver FACTS.md, secciones 3 y 4): IA aplicada y herramientas agénticas.
% El plugin de Claude Code que construyó para la plataforma que operaba enruta
% una petición en lenguaje natural a un script DETERMINISTA: no es
% entrenamiento de modelos, no es RAG, no es MLOps y no llama a ninguna API de
% LLM. La redacción nunca debe insinuar lo contrario.
%
% Las cifras del plugin (4 skills, 35 scripts en Python de unas 10,800 líneas,
% unas 3,200 líneas de SKILL.md) salen de la suite misma. El número de
% capacidades NO se publica: el índice autogenerado dice 17 y el frontmatter
% dice 18, y no hay una cifra única verificable. Los identificadores internos
% (boards, wikis, bastiones, orgs, nombres de agentes) se generalizan según la
% §4c de FACTS.md.
%
% El stack de desarrollo agéntico (Superpowers, Matt Pocock Skills, Ponytail,
% Karpathy Guidelines, Codebase Memory MCP) se publica como USADO, nunca como
% construido: las cinco herramientas son proyectos públicos de terceros (§3 de
% FACTS.md). El servidor MCP indexa el repositorio en un grafo de conocimiento
% persistente con análisis estático local y no lleva ningún modelo dentro: nunca
% debe describirse como almacén de embeddings, base vectorial o búsqueda
% semántica (§4 de FACTS.md).
%
% El stack es de la entrega MÓVIL. Lo que construyó en Veradat es otra cosa:
% skills propias para las tareas de la organización, especializadas en automatizar
% procesos DevOps. Los dos conjuntos no se mezclan (§4 de FACTS.md).
%
% Mantener cifras, empleadores y alcance alineados con lang/ai/en.tex y con
% los demás perfiles; FACTS.md es la fuente de verdad de cada afirmación.

\resumeaddress{Ciudad de México, México \textbar\ UTC-6 (US Central)}
\resumejobtitle{Ingeniero de IA \textbar\ Claude Code, skills de agentes y desarrollo agéntico}
\resumesummary{Ingeniero de IA con trayectoria en desarrollo de software desde 2017, enfocado en IA aplicada y herramientas agénticas sobre infraestructura en producción. He construido el \textbf{plugin de Claude Code} de la plataforma: \textbf{4 skills} que leen una petición en lenguaje natural, detectan la intención y la ejecutan con \textbf{35 scripts deterministas en Python}, documentados en cerca de \textbf{3{,}200 líneas de SKILL.md}. La unidad de trabajo es el contrato ejecutable entre la intención y el script. En su propio trabajo de entrega, el código se ha escrito con agentes de IA, sobre un flujo dirigido por specs. Antes, 20 pipelines de CI/CD para tres equipos y \textbf{21 clústeres de preproducción}.}
\resumesummaryheading{Sobre mí}
\makeheader

% Sidebar column. In the ATS layout \sidecol is a no-op and this prints inline.
\sidecol{%
\section*{Competencias Técnicas}
\skillrow{IA y herramientas agénticas}{Claude Code, SKILL.md, Python, Bash, YAML, JSON}
\skillrow{Plataformas y orquestación}{Kubernetes, OpenShift, Docker, Helm}
\skillrow{Nube e infraestructura como código}{AWS (EC2, S3, VPC, Security Groups, IAM, RDS, CloudWatch, Route~53), Azure, Terraform, Ansible, Linux}
\skillrow{CI/CD y entrega}{Azure DevOps, Jenkins, GitHub Actions, Nexus, GitHub, Maven, Microsoft Project}
\skillrow{Datos y mensajería}{Oracle, PostgreSQL, Kafka, RabbitMQ, IBM~MQ}
\skillrow{Desarrollo}{Java, Spring Cloud, JPA, APIs REST, TypeScript, React, React Native, Expo}
\skillrow{Observabilidad, seguridad y cumplimiento}{Elastic Stack (ELK), Splunk, ServiceNow, ISO 27001 (implementación y operación de controles), SonarQube Cloud, Snyk, Mend}
}

% Main column.
\maincol{%
\section*{Experiencia Profesional}

\resumeentry{Ingeniero DevOps}{Veradat Smart Network S.A. de C.V.}{Ago 2024 -- Oct 2026}{Ciudad de México}
\begin{itemize}
    \item Diseñé y mantuve \textbf{20 pipelines de CI/CD en Azure DevOps} para 3 equipos (frontend, backend y QA), reduciendo el tiempo de despliegue de 70 microservicios de \textbf{6 horas a 1 hora} (83\% menos tiempo). Desarrollé los charts de Helm que los tres equipos adoptaron como estándar de despliegue, y registré el trabajo de infraestructura en Microsoft Project por sprint.
    \item Construí \textbf{skills custom de Claude Code} para las tareas de la organización, especializadas en automatizar procesos DevOps: \textbf{4 skills} que leen una petición en lenguaje natural, detectan la intención y la ejecutan con un script determinista en Python -- \textbf{35 scripts, unas 10{,}800 líneas} -- documentados en unas \textbf{3{,}200 líneas de SKILL.md}, con un script que genera automáticamente el índice de capacidades. Las skills parsean el YAML y el JSON de la propia plataforma y ejecutan operaciones reales: validar el estado de releases y work items, verificar builds contra su rama, redesplegar un nodo con las versiones que ya corre, auditar la configuración del nodo contra el manual de la plataforma, leer secretos de aplicación en AWS Secrets Manager y reportar vulnerabilidades de Snyk.
    \item Revisé la arquitectura hexagonal del equipo, que desacopla cada microservicio de sus dependencias tecnológicas (PostgreSQL, Oracle, RabbitMQ, Kafka, IBM~MQ, SQS), y generé la imagen Docker para cada combinación tecnológica que cada microservicio requiere.
    \item Administro \textbf{21 clústeres de preproducción}: 9 en QA sobre Kubernetes, 4 en DEV y 4 en INTGR sobre OpenShift, 4 en Certificación sobre Kubernetes, con sondas readiness/liveness que controlan el tráfico en actualizaciones progresivas y disparan la recuperación, y monitoreo y alertas con Elastic Stack que redujeron el \textbf{tiempo de detección de incidentes de QA de 60 a 30 minutos} (50\% menos tiempo).
    \item Aprovisioné infraestructura en \textbf{AWS y Azure con Terraform} (EC2, S3, VPC, Security Groups, IAM, RDS, CloudWatch y Route~53), junto con máquinas virtuales, balanceadores de carga y gestión de secretos en ambas nubes, automaticé la configuración de 10 servidores Linux con Ansible, reduciendo el aprovisionamiento de \textbf{2 horas a 20 minutos} (83\% menos tiempo), y automatizé \textbf{20 procesos de infraestructura con shell scripting}. Realicé además backups y rollbacks de las bases de datos PostgreSQL y Oracle que respaldan los 70 microservicios, asegurando la recuperabilidad de los datos durante despliegues y respuesta a incidentes.
    \item Implementé y operé \textbf{controles ISO 27001} en gestión de secretos, control de acceso en nube y escaneo de dependencias, y remedié vulnerabilidades críticas y altas detectadas por SonarQube Cloud, Snyk y Mend mediante actualizaciones de paquetes del sistema operativo y dependencias Maven.
    \item Coordiné la interlocución técnica con stakeholders del cliente, incluido el reporte formal de seguridad que recibe su área de seguridad, y entre los 3 equipos (frontend, backend, QA), mediante dashboards y boards de Azure DevOps y una wiki compartida como canal asincrónico.
\end{itemize}

\resumeentry{Ingeniero Fullstack + DevOps}{IDS Comercial S.A. de C.V. \textbar\ Cliente: Banamex}{Ene 2021 -- Ago 2024}{México, remoto}
\begin{itemize}
    \item Desarrollé y mantuve \textbf{10 microservicios Java/Spring Cloud} para flujos de trabajo de ejecutivos bancarios que atendían aproximadamente \textbf{1{,}000 transacciones diarias por sucursal}, y los estandaricé con un POM padre multimódulo de Maven y gestión de dependencias compartida, garantizando un versionado consistente y reduciendo el esfuerzo de actualización en todos los servicios.
    \item Me apegaba a los \textbf{controles ISO 27001} en este entorno bancario durante toda mi colaboración, y construí pipelines de CI/CD con Jenkins para automatizar la compilación y el despliegue de la aplicación bancaria.
    \item Desarrollé funcionalidades frontend con React, Node.js y TypeScript, y administré bases de datos Oracle.
    \item Atendí y prioricé tickets de infraestructura en ServiceNow bajo un \textbf{SLA de 1 día}, revisando logs para troubleshooting y levantando incidentes al área correspondiente, y automatizé \textbf{10 procesos de infraestructura con shell scripting}. Monitoreé producción con Splunk y atendí incidentes del entorno de QA (tiempo promedio de resolución de 2 a 3 horas; casos complejos de hasta una semana).
\end{itemize}

\resumeentry{Ingeniero de Software Fullstack}{Proyecto continuo para Telcel, mediante GLOBAL HITSS y Everis}{Jul 2018 -- Dic 2020}{Ciudad de México}
\textit{GLOBAL HITSS: jul 2018 -- sep 2019; Everis: sep 2019 -- dic 2020.}\par\vspace{3pt}
\begin{itemize}
    \item Migré microservicios de un servidor de aplicaciones Tomcat a OpenShift, conteniendo la capa de integración construida con Java, Spring Cloud y JPA que traducía protocolos legados de telecomunicaciones a APIs REST, conectando sistemas existentes con servicios modernos.
    \item Implementé pipelines de CI/CD y contenericé el servicio con Docker sobre Kubernetes. Automaticé procesos operativos con Bash y gestioné repositorios de código y artefactos con GitHub y Nexus.
\end{itemize}

\section*{Proyectos Destacados e Independientes}

% mobile-begin
\resumeentry{Desarrollo de Aplicaciones Móviles (React Native)}{Proyectos freelance e independientes}{2023 -- 2026}{Ciudad de México}
\begin{itemize}
    \item Defino el contrato de entrega de apps para clientes en React Native y Expo, expresado como un artefacto ejecutable: \textbf{12 ADRs, 20 specs y 12 runbooks}, un \textbf{git worktree por slice} con su propia base de datos y un GitHub issue tracker numerado alimentado por commits convencionales.
    \item Establezco las puertas de release que el trabajo tiene que atravesar: \textbf{6 workflows de GitHub Actions} con jobs de prueba fragmentados sobre \textbf{95 archivos de prueba}, una \textbf{puerta de promoción} entre staging y producción, un \textbf{despliegue a producción con aprobación} y canales de EAS Update, sobre un monorepo con Bun workspaces, Turborepo y Biome.
    \item Construí esa entrega con agentes de IA sobre un flujo dirigido por specs: metodología \textbf{Superpowers} (spec, plan, TDD, subagentes, revisión), \textbf{Matt Pocock} Skills, \textbf{Ponytail}, \textbf{Karpathy} Guidelines y un grafo del repositorio con \textbf{Codebase Memory MCP}.
\end{itemize}
% mobile-end

\resumeentry{Sistema de Gestión de Cumplimiento de la Ley de Subcontratación}{GDP administración S de RL de CV (Grupo de la Parra)}{Jun 2021 -- Dic 2021}{Ciudad de México}
\begin{itemize}
    \item Lideré el diseño técnico y la arquitectura end-to-end, como líder técnico, de una plataforma en producción para gestionar el cumplimiento de la \textbf{Ley Federal del Trabajo en materia de subcontratación}, cubriendo una fuerza laboral subcontratada de \textbf{5{,}000 empleados} y múltiples empresas proveedoras y receptoras, y acompañé al equipo del cliente durante el diseño.
    \item Construí un \textbf{motor de auditoría de cumplimiento} que valida reglas regulatorias (existencia de CFDI por empleado vs. fechas de pago, pagos de IMSS/INFONAVIT/retenciones vs. calendario, alta ante la Secretaría del Trabajo) y emite \textbf{alertas automatizadas} por vencimientos de contratos y plazos de entrega de información, entregado con un \textbf{panel de BI de cumplimiento} (\textbf{SAT, IMSS, INFONAVIT, STPS}) y KPIs para las áreas legal y antifraude.
\end{itemize}
}

% Back to the sidebar.
\sidecol{%
\educationentry{Educación}{IPN -- Escuela Superior de C\'omputo (ESCOM)}{Ingeniería en Sistemas Computacionales}{2014 -- 2018}{Ciudad de México}

\section*{Certificación, Idiomas y Programación}
\skillrow{Certificación}{Java SE -- Oracle, Certificado Asociado (2013). \href{https://www.youracclaim.com/badges/583548b1-06fa-4dbf-889d-419e6c6c2ec1/public_url}{Credencial}; Scrum Fundamentals -- SCRUMstudy (2018).}
\skillrow{Idiomas}{Español nativo; inglés B2 (intermedio alto)}
\skillrow{Programación competitiva}{Participante, Google Code Jam (2019); equipo de selección de etapa de ACM-ICPC (2015--2017).}
}
